\documentclass[10pt,a4paper]{article}
\usepackage{amsmath,amssymb,amsthm,graphicx}
\usepackage{hyperref}
\usepackage{fontspec}
\title{Resolution lower bounds for functionally encoded Nisan-Wigderson generators via heavy width}
\author{Dmitry Sokolov}
\date{doi:10.4230/LIPIcs.CCC.2022.15; CC BY 4.0}
\begin{document}
\maketitle
\noindent\textit{Editorial scope.} Counted original Theorem14, the main Resolution lower bound for the functional NW encoding, with author expander and balanced-base-function hypotheses unchanged. Complete heavy-width, random self-reduction, local-consistency and AppendixA continuation proofs included; support Theorems23/26 and random-expander application15 are not separately counted. Mathematical assertions and parameter ranges are retained as published.\par\medskip
\noindent\textit{Allowed standard prerequisites (editorial).}\begin{itemize}
\item Elementary Boolean formula semantics and finite probability/union bounds. All source-local closure, expander peeling, self-reduction, balancedness and locally-consistent-assignment arguments retained, including deferred AppendixA graph lemmas. No main step delegated to cited background.
\end{itemize}\medskip
% Faithful official CCC2022 article15 mathematical transcription. Original author numbers retained.
\section*{2 Preliminaries}
Let $x$ be a propositional variable, i.e., a variable that ranges over the set $\{0,1\}$. A literal of $x$ is either $x$ (denoted sometimes as $x^1$) or $\neg x$ (denoted sometimes as $x^0$). A clause $C:=x_1^{c_1}\vee x_2^{c_2}\cdots\vee x_k^{c_k}$ is a disjunction of literals where $c_1,c_2,\ldots,c_k\in\{0,1\}$. A CNF formula $\varphi:=C_1\wedge\cdots\wedge C_m$ is a conjunction of clauses. We think of clauses and CNF formulas as sets: order is irrelevant and there are no repetitions.
A Resolution proof $\pi$ of an unsatisfiable CNF formula $\varphi$ is an ordered sequence of clauses $\pi:=C_1,\ldots,C_s$ such that $C_s=\emptyset$ is an empty clause and for each $i\in[s]$ either $C_i$ is a clause in $\varphi$ or there exist $j,k<i$ such that $C_i$ is derived from $C_j$ and $C_k$ by the resolution rule
\[\frac{C\vee x\qquad D\vee\neg x}{C\vee D}\]
or by the weakening rule
\[\frac CD\quad[C\subseteq D].\]
A partial assignment or a restriction on a function $f$ (or a formula $\varphi$) is a mapping $\rho:\operatorname{Vars}(f)\to\{0,1,*\}$. We let $\operatorname{supp}(\rho):=\rho^{-1}(\{0,1\})$ denote the set of assigned variables. The restriction of a function $f$ (or a formula $\varphi$) by $\rho$, denoted $f|_\rho$ ($\varphi|_\rho$), is the Boolean function (propositional formula) obtained from $f$ (from $\varphi$, respectively) by setting the value of each $x_i\in\operatorname{supp}(\rho)$ to $\rho(x_i)$ and leaving each $x_i\notin\operatorname{supp}(\rho)$ unassigned.
The size of a partial assignment $\rho$ is the size of the $\operatorname{supp}(\rho)$. We denote it by $|\rho|$.
\paragraph{Definition 3.} Let $f:\{0,1\}^n\to\{0,1\}$. We say that $f$ is $(\delta,k)$-balanced for some $0<\delta\leq1/2$ and $k\geq0$ iff for any $b\in\{0,1\}$ and any partial assignment $\rho$ of size at most $k$ the size of $(f|_\rho)^{-1}(b)$ is at least $\delta\cdot2^{n-|\rho|}$.
Some examples.
\begin{itemize}
\item $\operatorname{Parity}(x_1,\ldots,x_n)$ is $(1/2,n-1)$-balanced;
\item $\mathrm{IP}:=\sum_{i=1}^{n/2}x_iy_i\bmod2$ is $(1/4,n/2-1)$-balanced;
\item a random function is $(1/4,n-\sqrt n)$-balanced (see Lemma 30 for the calculations).
\end{itemize}
\subsection*{2.1 Expanders and Closure}
We use the following notation: $N_G(S)$ is the set of neighbours of the set of vertices $S$ in the graph $G$, $\partial_G(S)$ is the set of vertices $u$ that are connected with $S$ by exactly one edge. We omit the index $G$ if the graph is evident from the context.
\paragraph{Definition 4.} A bipartite graph $G:=(L,R,E)$ is an $(r,\Delta,c)$-expander if all vertices $u\in L$ have degree at most $\Delta$ and for all sets $S\subseteq L$, $|S|\leq r$, it holds that $|N(S)|\geq c\cdot|S|$. Similarly, $G:=(L,R,E)$ is an $(r,\Delta,c)$-boundary expander if all vertices $u\in L$ have degree at most $\Delta$ and for all sets $S\subseteq L$, $|S|\leq r$, it holds that $|\partial(S)|\geq c\cdot|S|$.
In this context, a simple but useful observation is that
\[|N(S)|\leq|\partial(S)|+\frac{\Delta|S|-|\partial(S)|}2=\frac{\Delta|S|+|\partial(S)|}2,\]
since all non-unique neighbours have at least two incident edges. This implies that for any $\varepsilon\leq1/2$ if a graph $G$ is an $(r,\Delta,(1-\varepsilon)\Delta)$-expander then it is also an $(r,\Delta,(1-2\varepsilon)\Delta)$-boundary expander.
The next proposition is well known in the literature. In this form it was used in [10].
\paragraph{Proposition 5.} If $G:=(L,R,E)$ is an $(r,\Delta,c)$-boundary expander then for any set $S\subseteq L$ of size $k\leq r$ there is an enumeration $v_1,v_2,\ldots,v_k\in S$ and a sequence $R_1,\ldots,R_k\subseteq N(S)$ such that:
\[R_i=N(v_i)\setminus\left(\bigcup_{j=1}^{i-1}N(v_j)\right);\qquad|R_i|\geq c.\]
In particular, there is a matching on the set $S$.
\paragraph{Proof.} We create this sequence in reversed order. Since $|S|\leq r$ it holds that $|\partial(S)|\geq c|S|$ and there is a vertex $v_k\in S$ such that $|\partial(S)\cap N(v_k)|\geq c$. Let $R_k:=|\partial(S)\cap N(v_k)|$, and repeat the process on $S\setminus\{v_k\}$.\hfill$\triangleleft$
Let $G:=(L,R,E)$ denote a bipartite graph. Consider a closure operation that seems to have originated in [3,2].
\paragraph{Definition 6.} For a vertex set $U\subseteq R$ we say that a set $S\subseteq L$ is $(U,r,\nu)$-contained if $|S|\leq r$ and $|\partial(S)\setminus U|<\nu|S|$. For any set $J\subseteq R$ let $\mathrm{Cl}_{r,\nu}(J)$ denote an arbitrary but fixed set of maximal size such that $\mathrm{Cl}_{r,\nu}(J)$ is $(J,r,\nu)$-contained. We say that $\mathrm{Cl}_{r,\nu}(J)$ is a closure of $J$.
Note that for any $J\subseteq R$ and any positive $r,\nu$ the empty set is $(J,r,\nu)$-contained and closure is well-defined.
\paragraph{Lemma 7.} Suppose that $G$ is an $(r,\Delta,c)$-boundary expander and that $J\subseteq R$ has size $|J|\leq\Delta r$. Then $|\mathrm{Cl}_{r,\nu}(J)|<|J|/(c-\nu)$.
\paragraph{Proof.} By definition we have that $|\partial(\mathrm{Cl}_{r,\nu}(J))\setminus J|<\nu|\mathrm{Cl}_{r,\nu}(J)|$. Since $|\mathrm{Cl}_{r,\nu}(J)|\leq r$ by definition, the expansion property of the graph guarantees that $c|\mathrm{Cl}_{r,\nu}(J)|-|J|\leq|\partial(\mathrm{Cl}_{r,\nu}(J))\setminus J|$. The conclusion follows.\hfill$\triangleleft$
Suppose $J\subseteq R$ is not too large. Then Lemma 7 shows that the closure of $J$ is not much larger. Thus, after removing the closure and its neighbourhood from the graph, we are still left with a decent expander. The following lemma makes this intuition precise.
\paragraph{Lemma 8.} Let $J\subseteq R$ be such that $|J|\leq\Delta r$ and $|\mathrm{Cl}_{r,\nu}(J)|\leq r/2$ and let $G':=G\setminus(\mathrm{Cl}_{r,\nu}(J)\cup J\cup N(\mathrm{Cl}_{r,\nu}(J)))$. Then any set $S$ of vertices from the left side of $G'$, with size $|S|\leq r/2$, satisfies that $|\partial_{G'}(S)|\geq\nu|S|$.
\paragraph{Proof.} Suppose the set $S\subseteq L(G')$ violates the boundary expansion guarantee. Observe that $\mathrm{Cl}_{r,\nu}(J)$ and $S$ are both sets of size at most $r/2$. Furthermore, the set $(\mathrm{Cl}_{r,\nu}(J)\cup S)$ is $(J,r,\nu)$-contained in the graph $G$. As $\mathrm{Cl}_{r,\nu}(J)$ is a $(J,r,\nu)$-contained set of maximal cardinality, this leads to a contradiction.\hfill$\triangleleft$
\section*{3 Nisan–Wigderson PRG and Its Encoding}
Let $G:=(L,R,E)$ be a bipartite graph such that $L:=\{v_1,v_2,\ldots,v_m\}$, $R:=\{u_1,u_2,\ldots,u_n\}$ and each vertex in $L$ has degree $\Delta$. Also for each vertex $v\in L$ we fix some arbitrary enumeration of its neighbours. We identify the right part of this graph with a set of boolean variables $\{x_1,x_2,\ldots,x_n\}$ and the left part with a set of output bits. Based on this identity we introduce a pseudorandom generator $\mathcal F_{G,f}:\{0,1\}^n\to\{0,1\}^m$ that is defined by the graph $G$ and a family of the base functions $f_1,f_2,\ldots,f_m:\{0,1\}^\Delta\to\{0,1\}$ in the natural way: the $j$-th bit of output is defined by $f_j(u_{i_1},u_{i_2},\ldots,u_{i_\Delta})$ (here we use enumeration of neighbours of the vertex $v_j$) where $u_{i_k}\in N(v_j)$ is a set of neighbours of the vertex $v_j\in L$. We also use a notation $\mathrm{Vars}_j:=N(v_j)$.
We want to encode the question about inversion of the function $\mathcal F_G$ as a propositional formula. Following the [2] and [1] we allow to use ``local'' extension variables.
\subsection*{3.1 Functional Encoding}
Let $\mathcal F_{G,f}:\{0,1\}^n\to\{0,1\}^m$ be a pseudorandom generator based on the graph $G$ and base functions $f_1,f_2,\ldots,f_m:\{0,1\}^\Delta\to\{0,1\}$. Let $b\in\{0,1\}^m$ be an arbitrary point. We say that a boolean function $g$ is local iff there is some $i\in[m]$ such that $g$ depends only on $\mathrm{Vars}_i$. Let $\mathcal G$ be a collection of local functions.
For each local function $g\in\mathcal G$ we introduce a variable $y_g$. And we write a CNF formula $\mathrm{PRG}_{G,f,b}$ on variables $y_g$ which consists of the following disjunctions:
\begin{itemize}
\item $(y_{g_1}^{c_1}\vee y_{g_2}^{c_2}\vee y_{g_3}^{c_3}\cdots\vee y_{g_\ell}^{c_\ell})$, for all tuples $g_1,g_2,\ldots,g_\ell$ where $\ell\leq2^{2^\Delta}$ and all $c_1,c_2,\ldots,c_\ell\in\{0,1\}$ such that there is $i\in[m]$:
\begin{itemize}
\item $g_j$ depends only on $\mathrm{Vars}_i$ for all $j\in[\ell]$;
\item any assignment $a\in\{0,1\}^\Delta$ that satisfy the equality $f_i(a)=b_i$ also satisfy the equality $g_k=c_k$ for at least one $k\in[\ell]$. In other words the equality $g_k=c_k$ semantically follows from the equality $f_i(a)=b_i$.
\end{itemize}
\end{itemize}
Note that, in particular, $\mathrm{PRG}_{G,f,b}$ contains the following constraints:
\begin{itemize}
\item $(y_{f_i}^{b_i})$, for all $i\in[m]$;
\item $(\neg y_s\vee y_g)$, $(\neg y_s\vee y_h)$ and $(y_s\vee\neg y_g\vee\neg y_h)$ for all local functions $s,g,h$ such that:
\begin{itemize}
\item $s=g\wedge h$;
\item there is $i\in[m]$ such that $s,g,h$ depends only on $\mathrm{Vars}_i$.
\end{itemize}
\end{itemize}
We omit indices of $\mathrm{PRG}_{G,f,b}$ if it is clear from the context. Following [2] we say that it is the functional encoding. The following observation is a straightforward corollary from the definition.
\paragraph{Remark 10 (Aleckhovich et al. [2]).} Formula $\mathrm{PRG}_{G,f,b}$ is unsatisfiable iff $b\notin\operatorname{Im}(\mathcal F_G)$.
\subsubsection*{3.1.1 Assignments, Restrictions and Intuition}
If we have some total assignment $\rho$ to $x$ variables of $\mathcal F$ (we call it $x$-assignment) it can define an assignment on $y$ variables in the natural way:
\begin{itemize}
\item $y_{x_i}\leftarrow\rho(x_i)$;
\item if $g(x)$ is a local function then $y_g\leftarrow g(x)|_\rho$.
\end{itemize}
We denote this assignment by $\rho_y$. The intuition behind the considered encoding and assignment is that a value of a variable $y_g$ on $\rho_y$ corresponds to a value of a function $g$ on $\rho$.
We can also define this extension (of $x$-assignments to $y$ variables) for partial assignments as well. If we have some partial assignment $\rho$ to $x$ variables, we define an assignment $\rho_y$ in the following way: $y_g|_{\rho_y}:=y_{g|_\rho}$. We say that these assignments are normal. In this paper we consider only normal assignments.
Consider a clause $D:=(y_{g_1}^{c_1}\vee y_{g_2}^{c_2}\vee\cdots\vee y_{g_\ell}^{c_\ell})$ in $y$ variables where $c_i\in\{0,1\}$. Note, that under normal assignments:
\[y_g^0\equiv y_{1-g};\qquad y_g\vee y_{g'}\equiv y_{g\vee g'}.\]
We can use these equalities and group variables of the clause $D$. Let $B_1,B_2,\ldots,B_m$ be a sequence of subsets (or bags) of literals $\{y_{g_1}^{c_1},y_{g_2}^{c_2},\ldots,y_{g_\ell}^{c_\ell}\}$ such that:
\begin{itemize}
\item for all $i\in[m],j\in[\ell]$ if $y_{g_j}^{c_j}\in B_i$ then $\operatorname{Vars}(g_j)\subseteq\mathrm{Vars}_i$;
\item for all $j\in[\ell]$ there is at least one $i\in[m]$ such that $y_{g_j}^{c_j}\in B_i$.
\end{itemize}
Note that we can rewrite a clause $D$ in the equivalent form under normal assignments $D\equiv(y_{h_1}\vee y_{h_2}\vee\cdots\vee y_{h_m})$ where
\[h_i(x):=\bigvee_{y_g^c\in B_i}(1\oplus c\oplus g(x)).\]
We may think about a clause $D$ as about the following disjunction:
\[F:=\bigvee_{i\in[m]}(h_i(x)=1),\]
and we say that $F$ is a functional form of the clause $D$. Denote by
\[F|_\rho:=\bigvee_{i\in[m]}(h_i(x)|_\rho=1),\]
where $\rho$ is an $x$-assignment.
\paragraph{Remark 11.} Functional form is not unique.
To illustrate Remark 11 consider the following situation for some function $g$ such that $y_g\in D$: $\operatorname{Vars}(g)\subseteq(\mathrm{Vars}_i\cap\mathrm{Vars}_{i'})$. In this case we can put the literal $y_g$ into a bag $B_i$, into a bag $B_{i'}$, or into both of these bags (also into none of them if we have an opportunity to put it into some third bag).
On the one hand Remark 11 provides additional problems if we want to work with functional forms of the clauses since we should care about all possible functional forms. On the other hand non-uniqueness allows to deal with assignments.
\paragraph{Lemma 12.} Let $F:=\bigvee_{i\in[m]}(h_i(x)=1)$ be a functional form of a clause $D$. If $\rho$ is a partial $x$-assignment then $F|_\rho$ is a functional form of a clause $D|_{\rho_y}$.
\paragraph{Proof.} Follows from definition of the functional form. Let $D:=(y_{g_1}^{c_1}\vee y_{g_2}^{c_2}\vee\cdots\vee y_{g_\ell}^{c_\ell})$, and $B_1,B_2,\ldots,B_m$ be a collection of bags that generates the functional form $(y_{h_1}\vee y_{h_2}\vee\cdots\vee y_{h_m})$.
Note that $D|_{\rho_y}=(y_{g_1}^{c_1}|_\rho\vee y_{g_2}^{c_2}|_\rho\vee\cdots\vee y_{g_\ell}^{c_\ell}|_\rho)$. We create a collection of bags $B'_i$ for a clause $D|_{\rho_y}$ and we put $y_{g_j}^{c_j}|_\rho\in B'_i$ iff $y_{g_j}^{c_j}\in B_i$. By construction it satisfy all required properties for bags and
\[\bigvee_{y_g^c\in B'_i}(1\oplus c\oplus g(x))=\bigvee_{y_g^c\in B_i}(1\oplus c\oplus g(x)|_\rho)=\left(\bigvee_{y_g^c\in B_i}(1\oplus c\oplus g(x))\right)|_\rho\]
that concludes the proof.\hfill$\triangleleft$
\paragraph{Remark 13.} The definition of functional form is ``syntactic'', or in other words it heavily depends on variables that appear in the clause $D$ and not only on the boolean function that is defined by it.
\section*{4 Lower Bound}
In this section we prove the main Theorem.

\paragraph{Theorem 14 (Formalization of Theorem 1).} Let $G:=(L,R,E)$ be an $(r,\Delta,(1-\varepsilon)\Delta)$-expander, where $|L|=m$, $|R|=n$. If $f_i$ is a family of $(1/4,3\varepsilon\Delta)$-balanced functions then for any $b\notin\operatorname{Im}(\mathcal F_{G,f})$ any resolution proof of $\mathrm{PRG}_{G,f,b}$ has size $\exp[\Omega(\varepsilon^5r^2/(2^{6\varepsilon\Delta}m))]$.

We defer the proof of this Theorem to section 4.4 and start with a plan of our proof.
We introduce an analog of the width measure on clauses with two important differences:
\begin{itemize}
\item we want to count number of output bits that are touched by a clause rather than number of variables;
\item we want to count only those outputs that we cannot erase from a clause ``for free''.
\end{itemize}
Let call this measure ``heavy width''.
We hit our proof by a random restriction. We will do it step by step and at each step we create an $x$-assignment $\sigma$ by choosing some output bit $v_i\in[m]$ and assign its neighbours in order to satisfy it. This assignment is equivalent to erasing some vertices from the right part of the graph. So after each step some output bit will not satisfy the expansion property of the graph. We say that these output bits are in danger and choose some $x$-assignment $\nu$ on its neighbours to satisfy them. We hit our proof by $(\sigma\cup\nu)_y$.
We repeat this process while we do not kill all clauses of big heavy width. At this step it is important that we deal with heavy width rather than usual width.
We prove a lower bound on the heavy width. For the sake of contradiction we assume that there is a proof of small heavy width. We trace a path in this resolution proof from the final clause to some axiom and maintain a partial assignment that:
\begin{itemize}
\item does not satisfy the current clause (note that this clause may have a large classical width, hence our assignments will not set this clause to a constant);
\item does not violate any axiom.
\end{itemize}
In the leaf these properties will give a contradiction.
\paragraph{Remark 16.} Note that if $f_i$ is a balanced function then $f_i(x)\oplus b_i$ is also a balanced function. Hence to simplify the notation wlog we assume that $b=0^n$ and we omit an index $b$ in the rest of the section. All the results holds for any $b\notin\operatorname{Im}(\mathcal F_{G,f})$.

\subsection*{4.1 The ``Heavy Width''}
In the classical restriction technique the notion of the width of a clause $C$ is used to estimate the probability that a random restriction will satisfy a clause. We give the next definition in order to save this property even if can deal with extension variables.
\paragraph{Definition 17.} Fix a formula $\mathrm{PRG}_{G,f}$. Let $C$ be a clause with functional form: $F:=\bigvee_{i=1}^m(h_i(x)=1)$. We say that $i$-th output bit is $\eta$-heavy in $F$ wrt $\mathrm{PRG}_{G,f}$ iff $\Pr_{z\leftarrow f_i^{-1}(0)}[h_i(z)=1]\geq\eta$. And the $\eta$-heavy width or $\mathrm{hw}_{\mathrm{PRG}_{G,f}}^\eta$ of $F$ is the number of $\eta$-heavy output bits in $F$.
This definition of width depends on the formula.
To justify this notion we may think about the ``information'' about $i$-th output bit in the clause $C$ (despite on the fact that it is defined for functional form). We pick a point $z\in\{0,1\}^\Delta$ that satisfy the constraint $f_i(z)=0$ uniformly at random. If the probability that we satisfy $C$ by this assignment is small then $C$ ``almost avoids'' $y$ variables that belongs to $i$-th output bit. In this case we pretend that the clause $C$ is independent of $i$-th output bit, otherwise the value of $C$ is heavily correlated with the value of $f_i$.
\paragraph{Remark 18.} The standard width measure can be considered as an $\eta$-heavy width measure. But in the different part of the proofs of classical resolution lower bounds we assume different parameters $\eta$.
\begin{itemize}
\item for the reduction from size to width: $\eta$ is an absolute positive constant (usually $1/2$);
\item for the width lower bound: $0<\eta<1/2^n$.
\end{itemize}
And it works since without extension variables for local functions we can state that: if an output bit is $\eta$-heavy for some $\eta>0$ then it also $\eta'$-heavy for some $\eta'\approx1/2$.
We define heavy width on functional form of the clause, that may give us potential problems due to the Remark 13.
\paragraph{Definition 19.} Let $\pi:=D_1,D_2,\ldots,D_s$ be a resolution proof of $\mathrm{PRG}_{G,f}$. And the $\eta$-heavy width $\mathrm{hw}_{\mathrm{PRG}_{G,f}}^\eta$ of the proof $\pi$ is the minimal natural number $w$ such that there is a sequence $F_1,F_2,\ldots,F_s$ where for all $i\in[s]$:
\begin{itemize}
\item $F_i$ is a functional form of $D_i$;
\item $\mathrm{hw}_{\mathrm{PRG}_{G,f}}^\eta$ of $F_i$ is at most $w$.
\end{itemize}
\subsection*{4.2 Size to Heavy Width Reduction}
In this section we present a random restriction argument that helps to reduce the question about size of proof to a question about $\eta$-heavy width of the proof for carefully chosen parameter $\eta$. Fix some $\mathrm{PRG}_{G,f}$.
Let define the key object that we use in our main Theorem.
\paragraph{Definition 20.} Let $G:=(L,R,E)$ be an $(r,\Delta,(1-\varepsilon)\Delta)$-expander. We say that an $x$-assignment $\rho$ of $\mathrm{PRG}_{G,f}$ is self-reduction iff there is a set $L_\rho\subseteq L$ such that:
\begin{itemize}
\item $|L_\rho|\leq\varepsilon^2r/16$;
\item $\rho$ assigns all and only variables from $N(L_\rho)$, moreover $\rho$ satisfy constraints from the set $L_\rho$;
\item $G\setminus(L_\rho\cup N(L_\rho))$ is an $(r,\Delta,(1-2\varepsilon)\Delta)$-expander.
\end{itemize}
The size of self-reduction is the size of the set $L_\rho$.
The next observation is trivial, but at the same it gives an opportunity to deal with heavy width measure since it is defined only for the PRG formulas.
\paragraph{Remark 21.} If $\rho$ is a self-reduction of $\mathrm{PRG}_{G,f}$ then $\mathrm{PRG}_{G,f}|_{\rho_y}$ is equivalent to $\mathrm{PRG}_{G',f'}$ under normal assignments where:
\begin{itemize}
\item $G':=G\setminus(L_\rho\cup N(L_\rho))$;
\item $f':=\{f'_1,\ldots,f'_m\}$ and $f'_i:=f_i|_\rho$.
\end{itemize}
We use the following algorithm to generate self-reductions.
\paragraph{Algorithm 1.} $r,\varepsilon$ are parameters.
\begin{enumerate}
\item $O_1:=\emptyset$ \hfill Set of active output bits
\item $G_1:=G$ \hfill $G_i=(L_i,R_i,E_i)$
\item $i:=1$
\item $\rho_1:=\emptyset$
\item For all $j\in[m]:p_j^1:=f_j$
\item while $i\leq\varepsilon^3r/32$ do
\item \quad Pick a vertex $v^i\in L_i$ uniformly at random
\item \quad Pick an $x$-assignment $\sigma_i\leftarrow(p_{v^i}^i)^{-1}(0)$ uniformly at random
\item \quad $O_{i+1}:=O_i\cup\{v^i\}$
\item \quad $G'_{i+1}:=G_i\setminus(\{v^i\}\cup N_{G_i}(v^i))$
\item \quad $B_i:=\arg\max\{|B|\mid B\subseteq L'_{i+1},|B|\leq r,|\partial_{G'_{i+1}}(B)|\leq(1-2\varepsilon)|B|\}$
\item \quad Pick an $x$-assignment $\nu_i$ on $N_{G'_{i+1}}(B_i)$ that satisfy all constraints from the set $B_i$
\item \quad $G_{i+1}:=G'_{i+1}\setminus(B_i\cup N_{G'_{i+1}}(B_i))$
\item \quad $\rho_{i+1}:=\rho_i\cup\sigma_i\cup\nu_i$
\item \quad For all $j\in[m]:p_j^{i+1}:=f_j|_{\rho_{i+1}}$
\item \quad $i:=i+1$
\end{enumerate}
return $\rho_i$
Following the Remark 21 note that $(\mathrm{PRG}_{G,f})|_{\rho_i^y}$ is equivalent to $\mathrm{PRG}_{G_i,p_i}$.
\paragraph{Lemma 22.} Let $G:=(L,R,E)$ be an $(r,\Delta,(1-\varepsilon)\Delta)$-expander graph such that $|L|=m$, $|R|=n$ and $10/\varepsilon\leq r$. If $f:=\{f_1,\ldots,f_m\}$ is a collection of $(1/4,3\varepsilon\Delta)$-balanced functions then Algorithm 1 generates a self-reduction of $\varphi:=\mathrm{PRG}_{G,f}$.
Before the proof we present an intuition about parameters. We are given a family of expander graphs that fixes some $\varepsilon$ (that we want to be as small as possible) and $\Delta$. We choose a family of balanced functions with proper parameters. Parameter $\varepsilon$ determines $\gamma:=2^{-\varepsilon\Delta}$, on the one hand it is just abbreviation, on the other hand it corresponds to the scaling factor $2^\rho$ from the definition of balanced function. In the classical Resolution lower bounds $\gamma$ is some constant (that is implicit and hidden inside the proof).
\paragraph{Proof.} Let $\ell:=\varepsilon^3r/32$ be the number of iterations of our algorithm.
By induction we show the following properties:
\begin{itemize}
\item $G_i$ is an $(r,\Delta,(1-2\varepsilon)\Delta)$-expander;
\item $|C_i|\leq\varepsilon^2r/32$,
\end{itemize}
where $C_i:=\bigcup_{j=1}^iB_j$. For proof see Proposition 29 in appendix A.
We have to show that on each iteration we can find some $x$-assignment $\nu_i$ that satisfy the requirements. Since $|C_i|\leq\varepsilon^2r/32$ that imply, in particular, that $|B_i|\leq\varepsilon^2r/32$.
Fix some iteration $i$. Since $G_i$ is an $(r,\Delta,(1-2\varepsilon)\Delta)$-expander then by Lemma 28 graph $G'_{i+1}$ is an $(r,\Delta,(1-3\varepsilon)\Delta)$-expander. Hence by Proposition 5 there is an enumeration $v^1,v^2,\ldots,v^{|B_i|}\in B_i$ and a sequence $R_1,\ldots,R_k\subseteq N_{G'_{i+1}}(S)$ such that for all $e\in[|B_i|]$:
\[R_e=N_{G'_{i+1}}(v^e)\setminus\left(\bigcup_{j=1}^{e-1}N_{G'_{i+1}}(v^j)\right);\qquad|R_e|\geq(1-3\varepsilon)\Delta.\]
We define the $x$-assignment $\nu_i$ step by step starting from $v^1$. Consider an auxiliary $x$-assignment $\kappa:=\rho_i\cup\sigma_i$. Since $f_{v^1}$ is a $(1/4,3\varepsilon\Delta)$-balanced function and $\kappa$ assigns at most $|N_G(v^1)|-|N_{G'_{i+1}}(v^1)|<3\varepsilon\Delta$ its variables then $f_{v^1}|_\kappa$ is not a constant and we define an $x$-assignment $\nu_i^{v^1}$ to $R_1$ variables to satisfy the constraint $f_{v^1}(x)=0$. We continue this process for vertices $v^j$ and $\kappa:=\rho_i\cup\sigma_i\cup\bigcup_{b=1}^{j-1}\nu_i^{v^b}$. The $x$-assignment $\nu_i:=\bigcup_{b=1}^{|B_i|}\nu_i^{v^b}$ satisfy all constraints from the set $B_i$ as desired.
At this moment we proved that we can realise all steps of our algorithm. The $x$-assignment $\rho_\ell$ assigns only variables from $N(O_\ell\cup C_\ell)$, hence $L_\rho:=O_\ell\cup C_\ell$. The first property of self-reduction is satisfied: $|L_\rho|\leq\ell+\varepsilon^2r/32\leq\varepsilon^2r/16$. The second follow from the construction. The third property was proved above. Moreover all intermediate $x$-assignments $\rho_i$ are also satisfy these properties.\hfill$\triangleleft$
\paragraph{Theorem 23.} Let $G:=(L,R,E)$ be an $(r,\Delta,(1-\varepsilon)\Delta)$-expander graph such that $|L|=m$, $|R|=n$ and $10/\varepsilon\leq r$. Fix $\gamma:=2^{-\varepsilon\Delta}$. Let $f:=\{f_1,\ldots,f_m\}$ be a collection of $(1/4,3\varepsilon\Delta)$-balanced functions, and $\pi:=D_1,\ldots,D_s$ be a resolution proof of $\varphi:=\mathrm{PRG}_{G,f}$.
If $s<\exp(\frac{\varepsilon^3r}{32}\cdot\frac13\gamma^6\frac wm)$ for some $w\in\mathbb N$ then there is a self-reduction $\rho$ such that $\mathrm{hw}_{\varphi|_{\rho_y}}^{\gamma^3}$ of $\pi|_{\rho_y}$ is at most $w$.
\paragraph{Proof.} Let $\theta:=F_1,F_2,\ldots,F_s$ where $F_i$ is a functional form of $D_i$.
We show that under the current assumptions Algorithm 1 whp give such an $x$-assignment. Let $\ell:=\varepsilon^3r/32$ be the number of iterations of our algorithm. By Lemma 22 it generates self-reduction $\rho$. We have to check that whp $\pi|_\rho$ is a proof of small heavy width. We do it for each line in the proof separately and moreover for all $i\in[s]$ we show that $(F_i)|_\rho$ has small heavy width.
One of the important difference of heavy width and a classical width that after application of a partial assignment it may increase since we also apply an assignment to the functions $f_i$ and change our formula. To avoid this problem we analyse $\mathrm{hw}_{\varphi}^{\gamma^6/3}$ rather than $\mathrm{hw}_{\varphi}^{\gamma^3}$ of the clauses and show that for any $F\in\theta$:
\begin{itemize}
\item $\mathrm{hw}_{\varphi|_{\rho_i^y}}^{\gamma^6/3}(F|_{\rho_i})$ is small for any $i\leq\ell$ than it cannot ``grow'' to much in the end (in terms of $\mathrm{hw}_{\varphi}^{\gamma^3}$);
\item if $\mathrm{hw}_{\varphi|_{\rho_i^y}}^{\gamma^6/3}(F|_{\rho_i})$ is big enough for some $i\leq\ell$ it will be killed with good enough probability on $i+1$-th iteration.
\end{itemize}
We start with the first part of the proof. Let $F\in\theta$ be a functional form of a clause from $\pi$ and $F$ is defined by the functions $h_1,h_2,\ldots,h_m$. Fix some $i\leq\ell$ and pick some alive output bit $v\in L\setminus L_\rho$. We remind a notation $p_v^i:=f_v|_{\rho_i}$. Output bit $v$ is alive and graph $G_\ell$ is an $(r,\Delta,(1-3\varepsilon)\Delta)$-expander, hence $x$-assignments $\rho_i$ and $\rho_\ell$ can assign at most $3\varepsilon\Delta$ variables from $N(v)$. Thus for all $i\leq\ell$:
\begin{align*}
\Pr_{z\leftarrow(p_v^\ell)^{-1}(0)}[h_v(z)=1]
&\leq\frac{|h_v^{-1}(1)\cap(p_v^\ell)^{-1}(0)|}{|(p_v^\ell)^{-1}(0)|}\\
&\leq\frac{|h_v^{-1}(1)\cap(p_v^i)^{-1}(0)|}{|(p_v^\ell)^{-1}(0)|}&&\rho_i\subseteq\rho_\ell\\
&=\frac{|h_v^{-1}(1)\cap(p_v^i)^{-1}(0)|}{|(p_v^i)^{-1}(0)|}\cdot\frac{|(p_v^i)^{-1}(0)|}{|(p_v^\ell)^{-1}(0)|}\\
&\leq\frac{|h_v^{-1}(1)\cap(p_v^i)^{-1}(0)|}{|(p_v^i)^{-1}(0)|}\cdot\frac{|f_v^{-1}(0)|}{|(p_v^\ell)^{-1}(0)|}&&p_v^i=f_v|_{\rho_i}\\
&\leq\frac{|h_v^{-1}(1)\cap(p_v^i)^{-1}(0)|}{|(p_v^i)^{-1}(0)|}\cdot\frac{\frac342^\Delta}{\frac142^{\Delta-3\varepsilon\Delta}}&&f\text{ is balanced}\\
&\leq\Pr_{z\leftarrow(p_v^i)^{-1}(0)}[h_v(z)=1]\cdot3\cdot2^{3\varepsilon\Delta}.
\end{align*}
Hence for all $F\in\theta$ and all $v\in L$ if $v$ is $\gamma^3$-heavy in $F|_{\rho_\ell}$ wrt $\varphi|_{\rho_y}$ then $v$ is $\gamma^6/3$-heavy in $F|_{\rho_i}$ wrt $\varphi|_{\rho_i^y}$ for all $i\leq\ell$. And $\mathrm{hw}_{\varphi|_{\rho_i^y}}^{\gamma^3}(F|_{\rho_\ell})\geq w$ imply that $\mathrm{hw}_{\varphi|_{\rho_i^y}}^{\gamma^6/3}(F|_{\rho_i})\geq w$ for all $i\in\ell$.
Consider a clause $D$ in $\pi|_{\rho_{i-1}^y}$ and its functional form $F\in\theta$. Denote by $H$ the event that $v^i$ is $\gamma^6/3$-heavy output bit in $F$ wrt $\mathrm{PRG}_{G,f}|_{\rho_{i-1}}$. Clause $D$ is killed by $\rho_i^y$ with probability at least:
\begin{align*}
\Pr_{v^i,\sigma_i}[D|_{\sigma_i}=1]&\geq\Pr_{v^i,\sigma_i}[h_{v^i}(x)|_{\sigma_i^y}=1]\\
&\geq\Pr_{v^i,\sigma_i}[H]\cdot\Pr_{v^i,\sigma_i}[(h_{v^i}|_{\sigma_i^y}=1)\mid H]\\
&\geq\frac{|\{v^i\in[m]\mid H\}|}{m}\cdot\frac{\gamma^6}3.
\end{align*}
Hence for the clause $D\in\pi$ there are two ways.
\begin{itemize}
\item At some moment $i\leq\ell$ the $\mathrm{hw}_{\varphi|_{\rho_i^y}}^{\gamma^6/3}(F|_{\rho_i})\leq w$. In this case $D$ is not interesting for us anymore, since as we proved above $\mathrm{hw}_{\varphi|_{\rho_y}}^{\gamma^3}(F|_\rho)\leq w$.
\item If $\mathrm{hw}_{\varphi|_{\rho_i^y}}^{\gamma^6/3}(F|_{\rho_i})\geq w$ then the probability that the clause $D|_{\rho_i}$ is survived on $i+1$-th iteration is at most:
\[\Pr[D|_{\rho_i}\text{ is survived on }i+1\text{-th iteration}]\leq1-\frac wm\frac{\gamma^6}3.\]
And hence:
\begin{align*}
\Pr[D\text{ is survived after }\ell\text{ iterations}]&\leq\prod_i\Pr[D|_{\rho_i}\text{ is survived on }i+1\text{-th iteration}]\\
&\leq\left(1-\frac wm\frac{\gamma^6}3\right)^\ell,
\end{align*}
since $\mathrm{hw}_{\varphi|_{\rho_i^y}}^{\gamma^6/3}(F|_{\rho_i})\geq w$.
\end{itemize}
To conclude the proof note that
\begin{align*}
\Pr[\mathrm{hw}_{\varphi|_{\rho_\ell^y}}^{\gamma^3}(F|_{\rho_\ell})>w]
&\leq\left(1-\frac wm\frac{\gamma^6}3\right)^\ell=\left(1-\frac wm\frac{\gamma^6}3\right)^{\varepsilon^3r/32}\\
&<\left(1-\frac wm\frac{\gamma^6}3\right)^{16\log s\,m/(w\gamma^3)}<1/s.
\end{align*}
By the union bound over all $F\in\theta$ we conclude that:
\[\Pr[\exists F\in\theta,\mathrm{hw}_{\varphi|_{\rho_\ell^y}}^{\gamma^3}(F|_{\rho_\ell})>w]<1.\]
Or in other words there is an $x$-assignment $\rho$ that satisfy all required properties.\hfill$\triangleleft$
\subsection*{4.3 Heavy Width Lower Bound}
For the sake of contradiction assume that we have a proof $\pi:=(D_1,\ldots,D_s)$ of small heavy width. Starting from $D_s$ we trace the path $p$ in the dag of $\pi$ to the initial clause. During this process we maintain a partial $x$-assignment $\sigma$ such that in the clause $D\in p$ for any small set $S$ of initial clauses the $x$-assignment $\sigma$ can be extended for an $x$-assignment $\kappa\supseteq\sigma$ such that $S$ is satisfied by $\kappa$, but $D$ does not. That give us a contradiction in a leaf where $D$ should be one of the initial clauses.
This assignment $\sigma$ will assign neighbours of heavy output bits of the generator and some extension (closure) to make sure that the remaining graph (after removing assigned variables) is an expander. Since the remainder is an expander (in particular we assign not so many neighbours of any alive output bit) then the existence of the assignment $\kappa$ will follow from the fact other output bit are not heavy that means that there are a lot of points that satisfy the constraint but not satisfy the clause.
In this section we assume that $G:=(L,R,E)$ be an $(r,\Delta,(1-\varepsilon)\Delta)$-expander and $\mathrm{PRG}_{G,f}$ is based on this graph an some functions $f_i$ (again we assume that zero point not in $\operatorname{Im}(\mathcal F_{G,f})$ and state our result for this point, but it holds for all $b\notin\operatorname{Im}(\mathcal F_{G,f})$). All clauses deal with variables of $\mathrm{PRG}_{G,f}$. We also fix an abbreviation $\gamma:=2^{-\varepsilon\Delta}$.
Let start with auxiliary objects and lemmas.
\paragraph{Definition 24.} Let $\rho$ be a self-reduction of $\mathrm{PRG}_{G,f}$. Let $C$ be a clause and $F$ its functional form, $I_\eta$ be a set of $\eta$-heavy output bits wrt $\mathrm{PRG}_{G,f}|_{\rho_y}$. We say that an output bit $v$ is $(\eta,\nu)$-dangerous for $F$ iff $v\in\mathrm{Cl}_{r,\nu}(N_G(I_\eta\cup L_\rho))$. Denote this set by $D_{F,\rho}^{\eta,\nu}$.
Note that this definition make sense only if graph $G$ is an expander. Also note that $I_\eta\cup L_\rho\subseteq\mathrm{Cl}_{r,\nu}(N_G(I_\eta\cup L_\rho))$. For fixed parameters $\eta$ and $\nu$ and a clause $C$ with functional form $F$ we also say that an $x$-assignment $\sigma\supseteq\rho$ is $(\eta,\nu)$-locally consistent iff:
\begin{itemize}
\item $\sigma^{-1}(\{0,1\})=N(D_{F,\rho}^{\eta,\nu})$;
\item $C|_{\sigma_y}\not\equiv1$;
\item $\sigma$ satisfy all constraints that correspond to $D_{F,\rho}^{\eta,\nu}$.
\end{itemize}
The following Lemma is the heart of the proof. It says that locally consistent assignments cannot violate any constraint from our formula.
\paragraph{Lemma 25.} Let $f:=\{f_1,\ldots,f_m\}$ be a collection of $(1/4,3\varepsilon\Delta)$-balanced functions, $\gamma<1/8$ and $\rho$ be a self-reduction of $\mathrm{PRG}_{G,f}$. If $C$ is a clause with functional form $F$ such that $\mathrm{hw}_{\mathrm{PRG}_{G,f}|_\rho}^{\gamma^3}(F)\leq\varepsilon r/8$ and $\sigma$ is a $(\gamma^3,(1-2\varepsilon\Delta))$-locally consistent assignment then for any $J\subseteq L$ such that $|J|\leq\varepsilon r/4$, there is an extension $\kappa\supseteq\sigma$ such that:
\begin{itemize}
\item $\kappa^{-1}(\{0,1\})\supseteq N(D_{F,\rho}^{\gamma^3,(1-2\varepsilon)\Delta})\cup N(J)$;
\item $C|_{\kappa_y}\not\equiv1$;
\item $\forall v\in J,f_v(x)|_\kappa=0$.
\end{itemize}
\paragraph{Proof.} Let $D:=D_{F,\rho}^{\gamma^3,(1-2\varepsilon)\Delta}$ and $F:=\bigvee_i(h_i(x)=1)$. Let $I:=\mathrm{Cl}_{r,(1-2\varepsilon)\Delta}(N_G(J)\cup N_G(D))\setminus D$. By Lemma 7 $|I|\leq3r/8$. By the definition of closure $I\supseteq J\setminus D$.
By Lemma 8 graph $G':=G\setminus(D\cup N_G(D))$ is an $(r/2,\Delta,(1-2\varepsilon)\Delta)$-expander. By Proposition 5 there is an enumeration $v^1,v^2,\ldots,v^{|I|}\in I$ and a sequence $R_1,\ldots,R_{|I|}\subseteq N_{G'}(S)$ such that:
\[R_i=N_{G'}(v^i)\setminus\left(\bigcup_{j=1}^{i-1}N_{G'}(v^j)\right);\qquad|R_i|\geq(1-2\varepsilon)\Delta.\]
We define a family of $x$-assignments $\nu_i$ and $\kappa_i:=\bigcup_{j=1}^i\nu_i\cup\sigma$ step by step, starting from $\nu_1$ in the following way:
\begin{itemize}
\item $\nu_i^{(-1)}(\{0,1\})=R_i$;
\item $f_{v^i}(x)|_{\kappa_i}=0$;
\item $C|_{\kappa_i^y}\not\equiv1$.
\end{itemize}
We have to show the existence of such $\nu_i$. Note that $|R_i|\geq(1-2\varepsilon)\Delta$ hence $\kappa_{i-1}$ can assign at most $2\varepsilon\Delta$ variables in $N_G(v^i)$. Since $f_{v^i}$ is a balanced function:
\[|(f_{v^i}|_{\kappa_{i-1}})^{-1}(0)|\geq\frac14\gamma^22^\Delta\geq\frac14\gamma^2|f_{v^i}^{-1}(0)|.\]
Output bit $v^i$ is not $\gamma^3$-heavy hence there are at most $\gamma^3|f_{v^i}^{-1}(0)|$ different $x$-assignments to $R_i$ that maps $h_{v^i}$ to $1$, assuming that $\gamma<1/8$ we can find an assignment that maps $h_{v^i}$ to $0$ and satisfy the constraint $f_{v^i}(x)=0$. We define $\kappa:=\kappa_{|I|}$.
It remains to check that $\kappa_y$ does not satisfy $C$, that does not immediately follow from the construction due to Remark 13. To show this fact we use an expansion of underlying graph. For the sake of contradiction assume that $\kappa_y$ maps some literal $y_g^c\in C$ to $1$ and this literal belongs to bag $B_v$. Consider three cases.
\begin{enumerate}
\item $\kappa$ assigns all variables from $N(v)$. In this case $h_v(x)|_\kappa$ is mapped to $1$ since $h_v$ is a disjunction of $(1\oplus c\oplus g)$ with some function. that contradicts with the choice of $\kappa$.
\item $\kappa$ assigns at most $2\varepsilon\Delta$ variables from $N(v)$. Note that
\begin{align*}
\Pr_{z\leftarrow f_v^{-1}(0)}[h_v(z)=1]&\geq\Pr_{z\leftarrow f_v^{-1}(0)}[z\text{ cons. with }\kappa]&&\kappa\text{ maps }y_g\text{ to }1\\
&\geq\Pr_{z\leftarrow\{0,1\}^\Delta}[z\text{ cons. with }\kappa\wedge f_v(z)=0]\\
&\geq\Pr_{z\leftarrow\{0,1\}^\Delta}[z\text{ cons. with }\kappa]\cdot\Pr_{z\leftarrow\{0,1\}^\Delta}[f_v(z)=0\mid z\text{ cons. with }\kappa]\\
&\geq\gamma^2\Pr_{z\leftarrow\{0,1\}^\Delta}[f_v(z)=0\mid z\text{ cons. with }\kappa]\\
&\geq\gamma^2\Pr_{z\leftarrow\{0,1\}^\Delta}[f_v(z)|_\kappa=0]\\
&\geq\gamma^2/4\geq\gamma^3.&&f_v\text{ is balanced}
\end{align*}
But in this case $v\in D$ by definition of $D$. Hence $\kappa$ should assign all variable in $N(v)$.
\item $\kappa$ assigns at least $2\varepsilon\Delta+1$ variables from $N(v)$ but not all of them. That contradicts with the fact that $\kappa$ assigns variables from $N(I\cap D)=N(\mathrm{Cl}_{r,(1-2\varepsilon)\Delta}(N_G(J)\cup N_G(D)))$ and Lemma 8.
\end{enumerate}\hfill$\triangleleft$
\paragraph{Theorem 26.} Let $\varepsilon<1/3$, $G:=(L,R,E)$ be an $(r,\Delta,(1-\varepsilon)\Delta)$-expander graph such that $|L|=m$, $|R|=n$. Fix $\gamma:=2^{-\varepsilon\Delta}<1/8$. If $f:=\{f_1,\ldots,f_m\}$ is a collection of $(1/4,3\varepsilon\Delta)$-balanced functions then $\mathrm{hw}_{\mathrm{PRG}_{G,f}|_{\rho_y}}^{\gamma^3}$ of any resolution proof of $\mathrm{PRG}_{G,f}|_{\rho_y}$ is at least $\varepsilon^2r/16$ where $\rho$ is self-reduction.
\paragraph{Proof.} For the sake of contradiction assume that $\pi:=(D_1,\ldots,D_s)$ is a resolution proof of $\mathrm{PRG}_{G,f}|_{\rho_y}$ of $\mathrm{hw}_{\mathrm{PRG}_{G,f}|_{\rho_y}}^{\gamma^3}$ at most $\varepsilon^2r/16$. Let $\theta:=(F_1,F_2,\ldots,F_s)$ be a sequence of functional forms that is witnessing heavy width of $\pi$. For a disjunction $F_i\in\theta$ we denote $\mathcal D_i:=D_{F_i,\rho}^{\gamma^3,(1-2\varepsilon)\Delta}$.
For the clause $D_s$ with functional form $F_s$ an $x$-assignment $\rho$ is locally consistent, since graph $(G\setminus(L_\rho\cup N(L_\rho))$ is an $(r,\Delta,(1-2\varepsilon)\Delta)$-expander. We want to show that the existence of a locally consistent assignment $\kappa_D$ for some clause $D\in\pi$ with functional form $F\in\theta$ imply the existence of a locally consistent assignment for at least one of its predecessors. In this case the can trace the path from $D_s$ to some initial clause $D_k\in\pi\cap\mathrm{PRG}_{G,f}|_{\rho_y}$ with functional form $F_k$ and show the existence of a locally consistent assignment $\kappa_{D_k}$ for this clause.
Suppose that $D_k:=(y_{g_1}^{c_1}\vee y_{g_2}^{c_2}\vee y_{g_3}^{c_3}\cdots\vee y_{g_\ell}^{c_\ell})$. By construction of $\mathrm{PRG}_{G,f}$ we can find an $i\in[m]$ such that for all $j\in[\ell]$:
\begin{itemize}
\item $g_j$ depends only on $\mathrm{Vars}_i$;
\item the equality $g_j(x)=c_j$ semantically follows from the equality $f_i(x)=0$.
\end{itemize}
But in this case $i$ is $1$-heavy for $D_k$ and for any $x$-assignment $\sigma$ the condition $f_i(x)|_\sigma=0$ imply that $D_k|_{\rho_y}\equiv1$. This fact contradicts with the definition of locally consistent assignment.
Suppose a locally consistent assignment $\kappa$ exists for a clause $D_i\in\pi$ with functional form $F_i\in\theta$ and $D_a,D_b$ are predecessors of $D_i$ in $\pi$ with functional forms $F_a,F_b$ respectively. Note, that $\mathrm{hw}_{\varphi}^{\gamma^3}$ of these functional forms are at most $\varepsilon^2r/16$, hence Lemma 7 together with the upper bound $|L_\rho|\leq\varepsilon^2r/16$ imply that the sizes of $\mathcal D_i,\mathcal D_a,\mathcal D_b$ are at most $\varepsilon r/16+\varepsilon r/16=\varepsilon r/8$. By Lemma 25 we have an extension $\sigma\supseteq\kappa$ on $N_G(\mathcal D_a\cup\mathcal D_b)$ that satisfy constraints of $\mathrm{PRG}_{G,f}$ that correspond to $\mathcal D_a\cup\mathcal D_b$ but do not satisfy $D_i$. And since $\sigma$ do not satisfy $D_i$ it also do not satisfy at least one of its predecessor, wlog it is $D_a$. And the $x$-assignment $\sigma\cap N_G(\mathcal D_a)$ is a locally consistent for $D_a$ and $F_a$ as desired.\hfill$\triangleleft$
\subsection*{4.4 Proof of Theorem 14}
For the sake of contradiction assume that $\pi:=D_1,D_2,\ldots,D_s$ is a resolution proof of $\mathrm{PRG}_{G,f}$ and $s\leq\exp[\delta\varepsilon^5r^2/(2^{6\varepsilon\Delta}m)]$ for some $\delta\leq10^{-4}$.
Fix $w:=\varepsilon^2r/20$ and $\gamma:=2^{-\varepsilon\Delta}$. Note that:
\begin{align*}
\exp\left(\frac{\varepsilon^3r}{32}\cdot\frac13\gamma^6\frac wm\right)
&=\exp\left(\frac{\varepsilon^3r}{32}\cdot\frac13\gamma^6\frac{\varepsilon^2r}{20\cdot m}\right)\\
&\geq\exp\left(\frac{\varepsilon^5}{2000}\cdot\gamma^6\frac{r^2}m\right)\\
&>\exp\left(\delta\varepsilon^5\cdot\gamma^6\frac{r^2}m\right)\geq s,
\end{align*}
hence we can apply Theorem 23 that gives a self-reduction $\rho$. We hit the proof $\pi$ by $\rho_y$ and the proof $\pi|_{\rho_y}$ is a proof of $\mathrm{PRG}_{G,f}|_{\rho_y}$. Moreover the $\mathrm{hw}_{\mathrm{PRG}_{G,f}|_{\rho_y}}^{\gamma^3}$ of $\pi|_{\rho_y}$ is at most $w$.
Since $\rho$ is a self-reduction then by Theorem 26 any proof of $\mathrm{PRG}_{G,f}|_{\rho_y}$ requires $\mathrm{hw}_{\mathrm{PRG}_{G,f}|_{\rho_y}}^{\gamma^3}$ at least $\varepsilon^2r/16>w$. Contradiction.

\section*{A Missed Lemmas (necessary author continuation)}
At first we show a simple auxiliary statement.
\paragraph{Lemma 27.} Suppose that $G:=(L,R,E)$ is an $(r,\Delta,c)$-boundary expander and that $J\subseteq R$ has size $|J|\leq\Delta r$. Then if $X\subseteq L$ has size $|X|\leq r$ and $|\partial(X)\setminus J|\leq\nu|X|$ then $X\leq|J|/(c-\nu)$.
\paragraph{Proof.} The expansion property of the graph guarantees that $c|X|-|J|\leq|\partial(X)\setminus J|$. The conclusion follows.\hfill$\triangleleft$
\paragraph{Lemma 28.} Let $G:=(L,R,E)$ be an $(r,\Delta,(1-\varepsilon)\Delta)$-boundary expander. Then $G\setminus(\{v\}\cup\{N(v)\})$ is an $(r-1,\Delta,(1-3\varepsilon/2)\Delta)$ expander where $v\in L$ is an arbitrary vertex.
\paragraph{Proof.} Fix some $v\in L$ and denote $G':=G\setminus(\{v\}\cup\{N(v)\})$.
Consider some set of $S\subseteq(L\setminus\{v\})$ of size at most $r-1$ and denote $H:=N(S)\cap N(v)$. Since $G$ is an expander:
\[|\partial_G(S\cup v)|=|\partial_{G'}(S)|+\Delta-|H|\geq(1-\varepsilon)\Delta(|S|+1)\]
\[|\partial_{G'}(S)|\geq(1-\varepsilon)\Delta|S|-\varepsilon\Delta+|H|.\]
But from the other point of view:
\[|\partial_{G'}(S)|\geq|\partial_G(S)|-|H|\geq(1-\varepsilon)\Delta|S|-|H|.\]
Altogether:
\begin{align*}
|\partial_{G'}(S)|&\geq(1-\varepsilon)\Delta|S|-\min(|H|,\varepsilon\Delta-|H|)\\
&\geq(1-\varepsilon)\Delta|S|-\frac\varepsilon2\Delta\geq(1-3\varepsilon/2)\Delta|S|.
\end{align*}\hfill$\triangleleft$
\paragraph{Proposition 29 (Analog of [24]).} For all $i\leq\ell$:
\begin{itemize}
\item $G_i$ is an $(r,\Delta,(1-2\varepsilon)\Delta)$-expander;
\item $|C_i|\leq|\varepsilon^2r/32|$.
\end{itemize}
\paragraph{Proof.} At first we prove the second claim $|C_i|\leq\varepsilon^2r/32$ by induction. $C_0$ is an empty set. Suppose that $|C_{i-1}|\leq\varepsilon^2r/32$. There are two steps in the proof:
\begin{itemize}
\item we show that $|B_i|\leq r/3$ that give us an opportunity to use expansion property for the set $C_i$;
\item we give a lower bound on size $\partial_G(C_i)$ by using expansion property and the upper bound by the choice of $B_i$ that together give us an upper bound on size of $C_i$.
\end{itemize}
Let start with the first step. $|\partial_G(B_i)\setminus N_G(C_{i-1}\cup O_{i+1})|\leq|\partial_{G'_{i+1}}(B_i)|\leq(1-2\varepsilon)|B_i|$. By definition $|B_i|\leq r$ and hence by Lemma 27 $|B_i|\leq|N_G(C_{i-1}\cup O_{i+1})|/(\varepsilon\Delta)\leq r/4+\varepsilon r/32\leq r/3$. That concludes the first step.
\begin{align*}
(1-\varepsilon)\Delta|C_i|&\leq|\partial_G(C_i)|&&\text{by expansion}\\
&\leq\left|\bigcup_{j=1}^i(\partial_G(B_j)\setminus N_G(C_{j-1}))\right|\\
&\leq\left|\bigcup_{j=1}^i(\partial_G(B_j)\setminus(N_G(C_{j-1})\cup N(O_{j+1})))\cup N(O_{j+1})\right|\\
&\leq\left|\bigcup_{j=1}^i\partial_{G'_{j+1}}(B_j)\cup N(O_{i+1})\right|&&\text{by the choice of }B_j\\
&\leq(1-2\varepsilon)\Delta\sum_{j=1}^i|B_j|+|N(O_{i+1})|\\
&\leq(1-2\varepsilon)\Delta|C_i|+|N(O_{i+1})|.
\end{align*}
And hence $|C_i|\leq|N(O_{i+1})|/(\varepsilon\Delta)\leq\varepsilon^2r/32$ as desired.
The first claim we prove by contradiction. Pick the minimal $i$ such that $G:=G_i$ is not an $(r,\Delta,(1-2\varepsilon))$-boundary expander and $S\subseteq L$ be a witness of it, i.e. $|S|\leq r$ and $|\partial_G(S)|\leq(1-2\varepsilon)|S|$. As in previous case $|\partial_G(S)\setminus(N_G(C_{i-1})\cup O_\ell)|\leq|\partial_G(S)|\leq(1-2\varepsilon)|S|$ hence by Lemma 27 $|S|\leq|N_G(C_{i-1})\cup O_\ell|/(\varepsilon\Delta)\leq r/2$.
Consider a set $S\cup B_{i-1}$ and note that size of it at most $r$. $\partial_{G'_i}(S\cup B_{i-1})\subseteq\partial_{G_i}(S)\cup\partial_{G'_i}(B_{i-1})$ by definition of $G_i$. This implies $|\partial_{G'_i}(S\cup B_{i-1})|\leq(1-2\varepsilon)\Delta|S|+(1-2\varepsilon)\Delta|B_{i-1}|=(1-2\varepsilon)\Delta|S\cup B_{i-1}|$. That contradicts with the choice of $B_{i-1}$.\hfill$\triangleleft$

\begin{thebibliography}{25}
\bibitem[1]{ref1} Michael Alekhnovich, Eli Ben-Sasson, Alexander A. Razborov, and Avi Wigderson. Space complexity in propositional calculus. SIAM J. Comput., 31(4):1184–1211, 2002. doi:10.1137/S0097539700366735.
\bibitem[2]{ref2} Michael Alekhnovich, Eli Ben-Sasson, Alexander A. Razborov, and Avi Wigderson. Pseudor- andom generators in propositional proof complexity. SIAM J. Comput., 34(1):67–88, 2004. doi:10.1137/S0097539701389944.
\bibitem[3]{ref3} Michael Alekhnovich and Alexander A. Razborov. Lower bounds for polynomial calculus: Non- binomial case. Proceedings of the Steklov Institute of Mathematics, 242:18–35, 2003. Available at http://people.cs.uchicago.edu/~razborov/files/misha.pdf.Preliminary version in FOCS ’01.
\bibitem[4]{ref4} Albert Atserias and Víctor Dalmau. A combinatorial characterization of resolution width. J. Comput. Syst. Sci., 74(3):323–334, 2008. doi:10.1016/j.jcss.2007.06.025.
\bibitem[5]{ref5} Paul Beame and Toniann Pitassi. Propositional proof complexity: Past, present and future. Bull. EATCS, 65:66–89, 1998.
\bibitem[6]{ref6} Eli Ben-Sasson and Avi Wigderson. Short proofs are narrow - resolution made simple. J. ACM, 48(2):149–169, 2001. doi:10.1145/375827.375835.
\bibitem[7]{ref7} Matthew Clegg, Jeff Edmonds, and Russell Impagliazzo. Using the groebner basis algorithm to find proofs of unsatisfiability. In Proceedings of the Twenty-Eighth Annual ACM Symposium on the Theory of Computing, Philadelphia, Pennsylvania, USA, May 22-24, 1996, pages 174–183, 1996. doi:10.1145/237814.237860.
\bibitem[8]{ref8} Stephen A. Cook and Robert A. Reckhow. The relative efficiency of propositional proof systems. J. Symb. Log., 44(1):36–50, 1979. doi:10.2307/2273702.
\bibitem[9]{ref9} Susanna F. de Rezende, Jakob Nordström, Kilian Risse, and Dmitry Sokolov. Exponential resolution lower bounds for weak pigeonhole principle and perfect matching formulas over sparse graphs. In Shubhangi Saraf, editor, 35th Computational Complexity Conference, CCC 2020, July 28-31, 2020, Saarbrücken, Germany (Virtual Conference), volume 169 of LIPIcs, pages 28:1–28:24. Schloss Dagstuhl - Leibniz-Zentrum für Informatik, 2020. doi:10.4230/LIPIcs.CCC.2020.28.
\bibitem[10]{ref10} Konstantinos Georgiou, Avner Magen, and Madhur Tulsiani. Optimal sherali-adams gaps from pairwise independence. In Irit Dinur, Klaus Jansen, Joseph Naor, and José Rolim, editors, Approximation, Randomization, and Combinatorial Optimization. Algorithms and Techniques, pages 125–139, Berlin, Heidelberg, 2009. Springer Berlin Heidelberg.
\bibitem[11]{ref11} Russell Impagliazzo, Pavel Pudlák, and Jirí Sgall. Lower bounds for the polynomial calculus and the gröbner basis algorithm. Computational Complexity, 8(2):127–144, 1999. doi:10.1007/s000370050024.
\bibitem[12]{ref12} Jan Krajíček. Bounded arithmetic, propositional logic, and complexity theory, volume 60 of Encyclopedia of mathematics and its applications. Cambridge University Press, 1995.
\bibitem[13]{ref13} Jan Krajíček. On the weak pigeonhole principle. Fundamenta Mathematicae, 170:123–140, January 2001. doi:10.4064/fm170-1-8.
\bibitem[14]{ref14} Jan Krajíček. Dual weak pigeonhole principle, pseudo-surjective functions, and provability of circuit lower bounds. J. Symb. Log., 69(1):265–286, 2004. doi:10.2178/jsl/1080938841.
\bibitem[15]{ref15} Jan Krajíček. Proof complexity generators. In Hajo Broersma, Stefan S. Dantchev, Matthew Johnson, and Stefan Szeider, editors, Algorithms and Complexity in Durham 2006 - Proceedings of the Second ACiD Workshop, 18-20 September 2006, Durham, UK, volume 7 of Texts in Algorithmics, page 3. King’s College, London, 2006.
\bibitem[16]{ref16} Noam Nisan and Avi Wigderson. Hardness vs randomness. J. Comput. Syst. Sci., 49(2):149–167, 1994. doi:10.1016/S0022-0000(05)80043-1.
\bibitem[17]{ref17} Pavel Pudlák. Proofs as games. Am. Math. Mon., 107(6):541–550, 2000. URL: http: //www.jstor.org/stable/2589349.
\bibitem[18]{ref18} Alexander A. Razborov. Bounded arithmetic and lower bounds in boolean complexity. In Peter Clote and Jeffrey B. Remmel, editors, Feasible Mathematics II, pages 344–386, Boston, MA, 1995. Birkhäuser Boston.
\bibitem[19]{ref19} Alexander A. Razborov. Lower bounds for propositional proofs and independence results in bounded arithmetic. In Friedhelm Meyer auf der Heide and Burkhard Monien, editors, Automata, Languages and Programming, 23rd International Colloquium, ICALP96, Paderborn, Germany, 8-12 July 1996, Proceedings, volume 1099 of Lecture Notes in Computer Science, pages 48–62. Springer, 1996. doi:10.1007/3-540-61440-0\_116.
\bibitem[20]{ref20} Alexander A. Razborov. Improved resolution lower bounds for the weak pigeonhole prin- ciple. Electron. Colloquium Comput. Complex., 8(55), 2001. URL: http://eccc.hpi-web.de/eccc-reports/2001/TR01-055/index.html.
\bibitem[21]{ref21} Alexander A. Razborov. Resolution lower bounds for the weak functional pigeonhole principle. Theor. Comput. Sci., 303(1):233–243, 2003. doi:10.1016/S0304-3975(02)00453-X.
\bibitem[22]{ref22} Alexander A. Razborov. Pseudorandom generators hard for k-dnf resolution and polynomial calculus resolution. Ann. of Math., 181:415–472, 2015. doi:10.4007/annals.2015.181.2.1.
\bibitem[23]{ref23} Nathan Segerlind, Samuel R. Buss, and Russell Impagliazzo. A switching lemma for small restrictions and lower bounds for k-dnf resolution. SIAM J. Comput., 33(5):1171–1200, 2004. doi:10.1137/S0097539703428555.
\bibitem[24]{ref24} Dmitry Sokolov. (Semi)Algebraic proofs over $\pm$1 variables. In Konstantin Makarychev, Yury Makarychev, Madhur Tulsiani, Gautam Kamath, and Julia Chuzhoy, editors, Proccedings of the 52nd Annual ACM SIGACT Symposium on Theory of Computing, STOC 2020, Chicago, IL, USA, June 22-26, 2020, pages 78–90. ACM, 2020. doi:10.1145/3357713.3384288.
\bibitem[25]{ref25} Andrew Chi-Chih Yao. Theory and applications of trapdoor functions (extended abstract). In 23rd Annual Symposium on Foundations of Computer Science, Chicago, Illinois, USA, 3-5 November 1982, pages 80–91. IEEE Computer Society, 1982. doi:10.1109/SFCS.1982.45.
\end{thebibliography}
\end{document}
